% ======================================================================
%  Soutenance M1 — Joseph MANELA — 19 août 2026, 9h
%  Compilation : pdflatex soutenance  (x2)
%  Un seul PDF : diapos principales (1-22) + annexes questions (23+)
% ======================================================================
\documentclass[10pt]{beamer}

% --- Langue & fontes --------------------------------------------------
\usepackage[utf8]{inputenc}
\usepackage[T1]{fontenc}
\usepackage[french]{babel}
\usepackage{lmodern}

% --- Maths & tableaux -------------------------------------------------
\usepackage{amsmath,amssymb}
\usepackage{booktabs}
\usepackage{tikz}
\usetikzlibrary{arrows.meta,positioning}
\usepackage{pgfplots}
\pgfplotsset{compat=1.18}

% --- Charte (couleur du mémoire) --------------------------------------
\definecolor{bleuunistra}{RGB}{0,61,110}
\definecolor{grisclair}{RGB}{242,242,242}
\definecolor{vertok}{RGB}{0,110,60}
\definecolor{rougeko}{RGB}{160,30,30}

\usetheme{default}
\usecolortheme{seahorse}
\setbeamercolor{structure}{fg=bleuunistra}
\setbeamercolor{frametitle}{bg=bleuunistra, fg=white}
\setbeamercolor{title}{fg=bleuunistra}
\setbeamercolor{block title}{bg=bleuunistra!15, fg=bleuunistra}
\setbeamercolor{block body}{bg=bleuunistra!5, fg=black}
\setbeamercolor{alerted text}{fg=rougeko}
\setbeamerfont{frametitle}{series=\bfseries}
\setbeamertemplate{navigation symbols}{}
\setbeamertemplate{footline}{%
  \hfill\scriptsize\insertframenumber/\inserttotalframenumber\hspace{2mm}\vspace{2mm}}
\setbeamertemplate{itemize item}{\color{bleuunistra}$\blacktriangleright$}
\setbeamertemplate{itemize subitem}{\color{bleuunistra}--}

\newcommand{\edge}{\textit{edge}}
\newcommand{\oui}{\textcolor{vertok}{\textbf{OUI}}}
\newcommand{\non}{\textcolor{rougeko}{non}}

\title{Juger une stratégie de bourse :\\
  un protocole statistique de validation\\ sur données dépendantes}
\author[J. Manela]{Joseph MANELA\\[2mm]
  {\small Encadrant : Nicolas Poulin}}
\institute{Master Mathématiques et applications --- Parcours Statistique\\
  Université de Strasbourg}
\date{19 août 2026}

\begin{document}

% ======================================================================
%  1. TITRE
% ======================================================================
\begin{frame}
  \titlepage
\end{frame}

% ======================================================================
%  2. PLAN
% ======================================================================
\begin{frame}{Plan}
  \begin{enumerate}\setlength{\itemsep}{5mm}\large
    \item Le contexte
    \item Les données : 11 stratégies, dont un témoin
    \item L'escalade de prudence, en trois actes
    \item Verdict, limites, perspectives
  \end{enumerate}
\end{frame}

% ======================================================================
%  3. CONTEXTE
% ======================================================================
\begin{frame}{Le contexte : de l'agent au protocole}
  \vspace{2mm}
  \begin{center}
  \resizebox{\linewidth}{!}{%
  \begin{tikzpicture}[
      box/.style={draw=bleuunistra, thick, rounded corners, align=center,
                  minimum height=10mm, fill=bleuunistra!5, font=\small},
      fl/.style={-{Stealth}, thick, bleuunistra}]
    \node[box] (agent) {Agent IA};
    \node[box, right=10mm of agent] (strat) {Stratégie de trading};
    \node[box, right=10mm of strat, fill=bleuunistra!20, thick]
      (proto) {\textbf{PROTOCOLE}\\ \scriptsize (ce mémoire)};
    \node[box, right=10mm of proto] (verdict) {avantage réel ?\\ oui / non / ?};
    \draw[fl] (agent) -- node[above,font=\scriptsize]{crée} (strat);
    \draw[fl] (strat) -- node[above,font=\scriptsize]{soumise} (proto);
    \draw[fl] (proto) -- (verdict);
  \end{tikzpicture}}
  \end{center}
  \vspace{6mm}
  \begin{itemize}
    \item Avant de trader : \structure{vrai pouvoir de prédiction, ou accident ?}
  \end{itemize}
\end{frame}

% ======================================================================
%  3bis. MARCHE / STRATEGIE / TRADE
% ======================================================================
\begin{frame}{Qu'est-ce qu'une stratégie de trading ?}
  \vspace{2mm}
  \begin{center}
  \begin{tikzpicture}[xscale=1.55,yscale=1.35]
    \draw[-{Stealth}] (0,0) -- (6.6,0) node[below left,font=\scriptsize]{temps};
    \draw[-{Stealth}] (0,0) -- (0,2.2) node[right,font=\scriptsize]{prix};
    \draw[thick,bleuunistra] plot[smooth,tension=0.8] coordinates
      {(0.2,1.0) (0.8,1.3) (1.4,0.9) (2.0,0.6) (2.6,0.95) (3.2,0.8)
       (3.8,1.25) (4.4,1.6) (5.0,1.4) (5.6,1.7) (6.2,1.5)};
    \filldraw[vertok] (2.0,0.6) circle (1.8pt)
      node[below left=1mm and 0mm,font=\scriptsize\bfseries,vertok]{achat};
    \filldraw[rougeko] (4.4,1.6) circle (1.8pt)
      node[above=2mm,font=\scriptsize\bfseries,rougeko]{vente};
    \draw[dashed,gray] (2.0,0.55) -- (2.0,0);
    \draw[dashed,gray] (4.4,1.55) -- (4.4,0);
    \draw[{Stealth}-{Stealth},gray] (2.0,0.22) --
      node[above=0.5mm,font=\scriptsize,gray]{détention} (4.4,0.22);
  \end{tikzpicture}
  \end{center}
  \vspace{3mm}
  \begin{itemize}\setlength{\itemsep}{3mm}
    \item le \textbf{marché} : des titres dont le prix fluctue en continu
    \item un \textbf{trade} : acheter, détenir, revendre $\to$ un rendement
    \item une \textbf{stratégie} : la règle qui décide \structure{quand} entrer
      et sortir --- avec les seules données \textbf{passées}
  \end{itemize}
\end{frame}

% ======================================================================
%  4. DONNEES
% ======================================================================
\begin{frame}{Le banc d'essai : 11 stratégies, 45\,008 trades}
  \begin{columns}[T]
    \begin{column}{0.54\textwidth}
      \scriptsize
      \begin{tabular}{@{}ll@{}}
        \toprule
        Stratégie & Fondée sur \\
        \midrule
        \texttt{rsi\_classic/strict/trend} & indicateur RSI \\
        \texttt{ma\_crossover} & moyennes mobiles \\
        \texttt{db\_bottom}, \texttt{dt\_top} & double creux / sommet \\
        \texttt{hs\_classic/inverse} & épaule-tête-épaule \\
        \texttt{sr\_breakout/breakdown} & cassures \\
        \midrule
        \texttt{oracle} (témoin) & \alert{connaît le futur} \\
        \bottomrule
      \end{tabular}
    \end{column}
    \begin{column}{0.42\textwidth}
      \small
      \begin{itemize}\setlength{\itemsep}{2mm}
        \item S\&P~500 : 501 titres
        \item 2010 $\to$ 2026
        \item unité $=$ le trade
        \item 978 à 8\,513 par stratégie
      \end{itemize}
    \end{column}
  \end{columns}
  \vspace{5mm}
  \begin{itemize}\setlength{\itemsep}{3mm}
    \item 10 règles classiques, \textbf{passé seulement}
    \item l'oracle triche (30 jours d'avance) : avantage \textbf{réel par
      construction}
  \end{itemize}
\end{frame}

% ======================================================================
%  5. PIEGE 1 : GAIN BRUT -> EDGE
% ======================================================================
\begin{frame}{Acte I --- « J'ai gagné, donc ça marche »}
  \begin{center}
    \structure{Marché haussier : $\sim97\,\%$ des titres montent.}\\
    Un gain positif ne prouve rien --- il faut battre \emph{le hasard}.
  \end{center}
  \vspace{4mm}
  Pour chaque trade : 200 trades fictifs \textbf{appariés}
  (même titre, durée, sens --- date tirée au sort)
  \[
    \edge \;=\; \underbrace{\text{rendement de la stratégie}}_{\texttt{return\_net}}
      \;-\; \underbrace{\text{rendement moyen du hasard}}_{\texttt{rand\_return}}
  \]
\end{frame}

% ======================================================================
%  6. STUDENT
% ======================================================================
\begin{frame}{Le test de référence : Student unilatéral}
  $\bar e$ : \edge{} moyen \textbf{observé} (connu)\\[1mm]
  $\mu$ : \edge{} moyen \textbf{théorique} (inconnu --- c'est lui qu'on juge)
  \vspace{3mm}
  \[
    \begin{cases}
      H_0 : \mu = 0\\
      H_1 : \mu > 0
    \end{cases}
    \qquad\quad
    t = \frac{\bar e}{s/\sqrt{n}} \;\underset{H_0}{\sim}\; t_{n-1}
  \]
  \vspace{2mm}
  \begin{itemize}\setlength{\itemsep}{3mm}
    \item $t$ $=$ nombre d'\textbf{erreurs-types} entre $\bar e$ et zéro
    \item p-value faible $\Rightarrow$ le hasard n'est plus crédible
    \item chez nous : \edge{} moyens de $-0{,}025$ à $+0{,}131$ ---
      significatifs ?
  \end{itemize}
  \vspace{4mm}
  \begin{block}{Deux conditions}
    \textbf{normalité} $\to$ TCL \qquad\qquad
    \textbf{indépendance} $\to$ actes II et III
  \end{block}
\end{frame}

% ======================================================================
%  7. NORMALITE : SHAPIRO -> TCL
% ======================================================================
\begin{frame}{Normalité : Shapiro--Wilk rejette\ldots\ le TCL répond}
  \begin{columns}[T]
    \begin{column}{0.42\textwidth}
      \footnotesize
      \begin{tabular}{@{}lrr@{}}
        \toprule
        Stratégie & $W$ & $p$ \\
        \midrule
        \texttt{ma\_crossover} & $0{,}50$ & $10^{-80}$ \\
        \texttt{oracle}        & $0{,}59$ & $10^{-76}$ \\
        \texttt{rsi\_classic}  & $0{,}87$ & $10^{-53}$ \\
        \texttt{dt\_top}       & $\alert{0{,}95}$ & $10^{-38}$ \\
        \bottomrule
      \end{tabular}
      \vspace{2mm}

      {\footnotesize $W = \dfrac{b^2}{\mathrm{SCE}} \simeq 1$ sous $H_0$
      (normalité)}
      \vspace{1.5mm}

      {\scriptsize $\mathrm{SCE}=\sum_i(x_i-\bar x)^2$ : dispersion
      observée\\[0.5mm]
      $b^2$ : dispersion attendue si les valeurs triées suivaient une
      normale}
    \end{column}
    \begin{column}{0.54\textwidth}
      \small
      $t$ ne dépend que de la \textbf{moyenne}. Or :
      \[
        \frac{\bar e - \mu}{\sigma/\sqrt{n}}
        \;\xrightarrow[n\to\infty]{}\; \mathcal N(0,1)
      \]
      \begin{itemize}
        \item quelle que soit la loi de départ
      \end{itemize}
    \end{column}
  \end{columns}
\end{frame}

% ======================================================================
%  8. BONFERRONI
% ======================================================================
\begin{frame}{Onze tests, pas un : la correction de Bonferroni}
  \[
    \mathrm P(\text{au moins un faux positif})
      = 1 - 0{,}95^{11} \approx \alert{0{,}43}
  \]
  \vspace{2mm}
  \begin{block}{On inverse : risque global fixé à 5\,\% (inégalité de Boole)}
    \[
      \mathrm P\Bigl(\bigcup_i A_i\Bigr) \leq \sum_i \mathrm P(A_i)
      \qquad\Longrightarrow\qquad
      \alpha' = \frac{0{,}05}{11} \approx \mathbf{0{,}0045}
    \]
  \end{block}
\end{frame}

% ======================================================================
%  9. VERDICT ETAPE 1
% ======================================================================
\begin{frame}{Premier verdict (Student naïf)}
  \begin{center}
  \small
  \begin{tabular}{@{}lrrrc@{}}
    \toprule
    Stratégie & $n$ & \edge{} moyen & $p$ & Verdict \\
    \midrule
    \texttt{oracle} (témoin) & 6\,798 & $+0{,}0658$ & $\approx 0$ & \oui \\
    \midrule
    \texttt{rsi\_classic} & 7\,117 & $+0{,}0123$ & $1{,}2\times10^{-8}$  & \oui \\
    \texttt{rsi\_strict}  &   978  & $+0{,}1313$ & $5{,}5\times10^{-10}$ & \oui \\
    \texttt{rsi\_trend}   & 2\,789 & $+0{,}0021$ & $0{,}25$ & \non \\
    \texttt{db\_bottom}   & 8\,513 & $-0{,}0025$ & $0{,}98$ & \non \\
    \texttt{ma\_crossover}& 5\,417 & $-0{,}0248$ & $\approx 1$ & \non \\
    \multicolumn{5}{c}{\scriptsize autres : $p > 0{,}45$, \non} \\
    \bottomrule
  \end{tabular}
  \end{center}
  \vspace{4mm}
  \begin{center}
    Huit écartées, \structure{deux survivantes} $+$ le témoin.\\[2mm]
    \alert{Mais l'indépendance des trades n'a jamais été vérifiée.}
  \end{center}
\end{frame}

% ======================================================================
%  10. TROIS CANAUX
% ======================================================================
\begin{frame}{Acte II --- canal 1 : la dépendance de titre}
  \begin{columns}[T]
    \begin{column}{0.60\textwidth}
      \small
      Deux trades d'une même action subissent les mêmes chocs
      $\to$ les titres forment des \textbf{grappes}.
      \[
        \hat\rho = \frac{\mathrm{MSB}-\mathrm{MSW}}
                        {\mathrm{MSB}+(n_0-1)\,\mathrm{MSW}}
      \]
      \begin{itemize}\setlength{\itemsep}{2.5mm}
        \item MSB : variation \textbf{entre} les titres
        \item MSW : variation \textbf{dans} un titre
        \item trades d'un titre semblables $\Rightarrow$ MSW s'effondre
          $\Rightarrow$ $\hat\rho$ monte
      \end{itemize}
    \end{column}
    \begin{column}{0.36\textwidth}
      \scriptsize
      \begin{tabular}{@{}lr@{}}
        \toprule
        & $\hat\rho$ \\
        \midrule
        \texttt{oracle}        & $0{,}053$ \\
        \midrule
        \texttt{db\_bottom}    & $0{,}000$ \\
        \texttt{dt\_top}       & $0{,}000$ \\
        \texttt{hs\_classic}   & $0{,}087$ \\
        \texttt{hs\_inverse}   & $0{,}001$ \\
        \texttt{ma\_crossover} & $0{,}003$ \\
        \texttt{rsi\_classic}  & $0{,}000$ \\
        \texttt{rsi\_strict}   & $\alert{0{,}134}$ \\
        \texttt{rsi\_trend}    & $0{,}010$ \\
        \texttt{sr\_breakdown} & $0{,}051$ \\
        \texttt{sr\_breakout}  & $0{,}004$ \\
        \bottomrule
      \end{tabular}
    \end{column}
  \end{columns}
\end{frame}

% ======================================================================
%  10bis. ACP DU MARCHE
% ======================================================================
\begin{frame}{Canal 2 : la dépendance de marché}
  \begin{columns}[T]
    \begin{column}{0.47\textwidth}
      \centering
      \begin{tikzpicture}
      \begin{axis}[
        width=\linewidth, height=5.2cm,
        ybar, bar width=5pt,
        xlabel={Composante}, ylabel={Variance expliquée (\%)},
        xlabel style={font=\scriptsize}, ylabel style={font=\scriptsize},
        tick label style={font=\tiny},
        xmin=0.3, xmax=11.7, ymin=0, ymax=80,
        xtick={1,2,3,4,5,6,7,8,9,10,11},
        ytick={0,20,40,60,80},
        ymajorgrids=true, grid style={dotted,gray!40},
        axis lines=left,
      ]
      \addplot[fill=bleuunistra!75, draw=bleuunistra] coordinates {
        (1,73.0) (2,7.9) (3,5.1) (4,3.0) (5,2.3) (6,2.1)
        (7,1.9) (8,1.5) (9,1.3) (10,1.2) (11,0.6)
      };
      \end{axis}
      \end{tikzpicture}
    \end{column}
    \begin{column}{0.51\textwidth}
      \centering
      \begin{tikzpicture}
      \begin{axis}[
        width=\linewidth, height=5.2cm,
        xlabel={Axe 1 (73~\%)}, ylabel={Axe 2 (7,9~\%)},
        xlabel style={font=\scriptsize}, ylabel style={font=\scriptsize},
        tick label style={font=\tiny},
        xmin=0, xmax=0.42, ymin=-0.42, ymax=0.80,
        xtick={0,0.1,0.2,0.3,0.4},
        ytick={-0.4,-0.2,0,0.2,0.4,0.6,0.8},
        grid=both, grid style={dotted,gray!40},
        axis lines=left,
      ]
      \addplot[only marks, mark=*, bleuunistra, mark size=1.6pt] coordinates {
        (0.277,0.433) (0.307,0.023) (0.299,0.316) (0.244,0.674)
      };
      \addplot[only marks, mark=triangle*, black!65, mark size=2.1pt] coordinates {
        (0.310,-0.183) (0.321,-0.185) (0.255,-0.227) (0.328,-0.109)
        (0.335,-0.129) (0.304,-0.303) (0.321,-0.128)
      };
      \draw[dashed,gray!70] (axis cs:0,0) -- (axis cs:0.42,0);
      \node[font=\tiny,anchor=east,black!75] at (axis cs:0.244,0.674) {Serv.\,coll.~};
      \node[font=\tiny,anchor=east,black!75] at (axis cs:0.277,0.433) {Cons.\,base~};
      \node[font=\tiny,anchor=west,black!75] at (axis cs:0.299,0.316) {~Immobilier};
      \node[font=\tiny,anchor=west,black!75] at (axis cs:0.307,0.023) {~Santé};
      \node[font=\tiny,anchor=east,black!75] at (axis cs:0.328,-0.109) {Finance~};
      \node[font=\tiny,anchor=west,black!75] at (axis cs:0.335,-0.129) {~Industrie};
      \node[font=\tiny,anchor=east,black!75] at (axis cs:0.321,-0.128) {Matér.~};
      \node[font=\tiny,anchor=east,black!75] at (axis cs:0.310,-0.183) {Com.~};
      \node[font=\tiny,anchor=west,black!75] at (axis cs:0.321,-0.185) {~Cons.\,disc.};
      \node[font=\tiny,anchor=east,black!75] at (axis cs:0.255,-0.227) {Énergie~};
      \node[font=\tiny,anchor=west,black!75] at (axis cs:0.304,-0.303) {~Techno.};
      \end{axis}
      \end{tikzpicture}
    \end{column}
  \end{columns}
  \vspace{3mm}
  \begin{itemize}\setlength{\itemsep}{2mm}
    \item corrélation moyenne entre les onze secteurs : $0{,}70$
    \item un axe capte \alert{$73\,\%$} de la variance : le marché bouge d'un
      bloc
    \item axe 2 : défensifs ($\bullet$) contre cycliques
      ($\blacktriangle$)
  \end{itemize}
\end{frame}

% ======================================================================
%  10ter. CANAL 3
% ======================================================================
\begin{frame}{Canal 3 : la dépendance temporelle}
  \begin{columns}[T]
    \begin{column}{0.60\textwidth}
      \small
      Les trades \emph{durent} ($78\,\% > 1$ mois) : les détentions se
      recouvrent, les mois voisins se ressemblent.
      \[
        \hat\rho(1) = \frac{\sum_t (x_t-\bar x)(x_{t+1}-\bar x)}
                           {\sum_t (x_t-\bar x)^2}
      \]
      \begin{itemize}\setlength{\itemsep}{2.5mm}
        \item $x_t$ : \edge{} moyen des trades entrés le mois $t$
        \item fort pour les trades longs, faible pour les courts
      \end{itemize}
    \end{column}
    \begin{column}{0.36\textwidth}
      \scriptsize
      \begin{tabular}{@{}lr@{}}
        \toprule
        & $\hat\rho(1)$ \\
        \midrule
        \texttt{oracle}        & $0{,}18$ \\
        \midrule
        \texttt{db\_bottom}    & $0{,}36$ \\
        \texttt{dt\_top}       & $0{,}38$ \\
        \texttt{hs\_classic}   & $0{,}15$ \\
        \texttt{hs\_inverse}   & $0{,}31$ \\
        \texttt{ma\_crossover} & $\alert{0{,}56}$ \\
        \texttt{rsi\_classic}  & $0{,}23$ \\
        \texttt{rsi\_strict}   & $0{,}03$ \\
        \texttt{rsi\_trend}    & $0{,}20$ \\
        \texttt{sr\_breakdown} & $0{,}11$ \\
        \texttt{sr\_breakout}  & $0{,}05$ \\
        \bottomrule
      \end{tabular}
    \end{column}
  \end{columns}
\end{frame}

% ======================================================================
%  11. DEFF + BOOTSTRAP
% ======================================================================
\begin{frame}{Porte de l'action (1/2) : l'effet de sondage (DEFF)}
  \textbf{Définition} : le rapport entre la variance réelle de la moyenne
  (trades groupés par titres) et celle d'un sondage aléatoire simple.
  \[
    \mathrm{DEFF}
      = \frac{\mathrm{Var}(\bar x)_{\text{grappes}}}
             {\mathrm{Var}_{\mathrm{SAS}}(\bar x)}
      = 1 + (\bar m - 1)\,\hat\rho
  \]
  \[
    n_{\mathrm{eff}} = \frac{n}{\mathrm{DEFF}}
    \qquad\quad
    \widehat{\mathrm{se}}_{\mathrm{corr}} = \frac{s}{\sqrt n}\sqrt{\mathrm{DEFF}}
  \]
  {\footnotesize $\bar m$ : nombre moyen de trades par titre}
  \vspace{3mm}
  \begin{center}
  \footnotesize
  \begin{tabular}{@{}lrrc@{}}
    \toprule
    & $\hat\rho$ & DEFF & $p$ corrigée \\
    \midrule
    \texttt{oracle}       & $0{,}053$ & $1{,}67$ & $\approx 0$ \\
    \texttt{rsi\_classic} & $0{,}000$ & $1{,}00$ & $1{,}2\times10^{-8}$ \\
    \texttt{rsi\_strict}  & $0{,}134$ & $1{,}15$ & $7\times10^{-9}$ \\
    \bottomrule
  \end{tabular}
  \end{center}
\end{frame}

% ======================================================================
%  11bis. BOOTSTRAP
% ======================================================================
\begin{frame}{Porte de l'action (2/2) : le bootstrap par grappes}
  \begin{center}
  \resizebox{0.95\linewidth}{!}{%
  \begin{tikzpicture}[
      grap/.style={draw=bleuunistra, rounded corners, fill=bleuunistra!5,
                   font=\scriptsize, minimum width=17mm, minimum height=6mm},
      fl/.style={-{Stealth}, thick, bleuunistra}]
    \node[grap] (a) {titre A};
    \node[grap, below=1.5mm of a] (b) {titre B};
    \node[grap, below=1.5mm of b] (c) {titre C};
    \node[below=0mm of c, font=\scriptsize] (dots) {$\vdots$};
    \draw[fl] (1.6,-0.75) -- node[above,font=\scriptsize,align=center]
      {tirage avec remise\\ de titres \textbf{entiers}\\ $\times\,4\,000$}
      (4.2,-0.75);
    \begin{scope}[shift={(6.9,-1.75)}]
      \fill[rougeko!25] plot[domain=-2.4:-1.5,samples=25]
        (\x,{2.0*exp(-\x*\x/1.1)}) -- (-1.5,0) -- (-2.4,0) -- cycle;
      \draw[thick,bleuunistra] plot[domain=-2.5:2.5,samples=60]
        (\x,{2.0*exp(-\x*\x/1.1)});
      \draw[-{Stealth}] (-2.7,0) -- (2.8,0);
      \draw[dashed,rougeko] (-1.5,0) -- (-1.5,1.1)
        node[above,font=\scriptsize,rougeko] {$0$};
      \draw[dashed,gray] (0,0) -- (0,2.05);
      \node[above,font=\scriptsize,gray] at (0,2.05) {$\bar e$};
      \node[font=\scriptsize,rougeko,anchor=north,align=center]
        at (-1.8,-0.15) {p-value $=$ part des\\ tirages $\le 0$};
    \end{scope}
  \end{tikzpicture}}
  \end{center}
  \vspace{2mm}
  \begin{itemize}\setlength{\itemsep}{2.5mm}
    \item un titre tiré emporte \textbf{tous ses trades} : la dépendance
      interne voyage intacte dans chaque réplique
    \item aucune hypothèse de forme ; chez nous :
      $p < 2{,}5\times10^{-4}$ pour les trois
  \end{itemize}
\end{frame}

% ======================================================================
%  12. LES DEUX PORTES (methode)
% ======================================================================
\begin{frame}{Les deux portes : une unité réelle, un vote}
  \vspace{2mm}
  \begin{center}
  \begin{tikzpicture}[
      box/.style={draw=bleuunistra, thick, rounded corners, align=center,
                  minimum height=9mm, fill=bleuunistra!5, font=\scriptsize},
      fl/.style={-{Stealth}, thick, bleuunistra}]
    \node[box] (trades) {45\,008 trades\\ (dépendants)};
    \node[box, above right=2mm and 14mm of trades] (titres)
      {\textbf{Porte action}\\ \edge{} moyen \emph{par titre} : 431--501 votes};
    \node[box, below right=2mm and 14mm of trades] (mois)
      {\textbf{Porte période}\\ \edge{} moyen \emph{par mois} : 138--197 votes};
    \node[box, below right=2mm and 14mm of titres, fill=bleuunistra!20]
      (et) {\textbf{ET}\\ seuil Bonferroni\\ à chaque porte};
    \draw[fl] (trades) -- (titres);
    \draw[fl] (trades) -- (mois);
    \draw[fl] (titres) -- (et);
    \draw[fl] (mois) -- (et);
  \end{tikzpicture}
  \end{center}
  \vspace{4mm}
  \begin{center}
  \small
  \begin{tabular}{@{}lrrrrc@{}}
    \toprule
    & \multicolumn{2}{c}{Porte action} & \multicolumn{2}{c}{Porte période} & \\
    \cmidrule(lr){2-3}\cmidrule(lr){4-5}
    Stratégie & $k$ & $p_A$ & $k$ & $p_B$ & Verdict \\
    \midrule
    \texttt{oracle} (témoin) & 501 & $\approx 0$ & 193 & $\approx 0$ & \oui \\
    \midrule
    \texttt{rsi\_classic} & 501 & $\alert{0{,}059}$ & 195 & $\alert{0{,}039}$ & \non \\
    \texttt{rsi\_strict}  & 456 & $3{,}9\times10^{-4}$ & 138 & $\alert{0{,}26}$ & \non \\
    \bottomrule
  \end{tabular}
  \end{center}
  \vspace{2mm}
  \begin{itemize}
    \item perte d'information assumée : faux négatif plutôt que faux positif
  \end{itemize}
\end{frame}

% ======================================================================
%  14. ACTE III : LA CHAINE
% ======================================================================
\begin{frame}{Acte III --- la porte de la période tient-elle ?}
  Les mois voisins restent liés (recouvrement) : la suite des \edge{} mensuels
  est une \textbf{série temporelle}.
  \vspace{3mm}
  \begin{center}
  \begin{tikzpicture}[
      box/.style={draw=bleuunistra, thick, rounded corners, align=center,
                  minimum height=9mm, minimum width=30mm, fill=bleuunistra!5,
                  font=\scriptsize},
      fl/.style={-{Stealth}, thick, bleuunistra}]
    \node[box] (bp) {\textbf{1. Box--Pierce}\\ dépendance ?};
    \node[box, right=8mm of bp] (st) {\textbf{2. DF $+$ KPSS}\\ stationnarité ?};
    \node[box, right=8mm of st] (pacf) {\textbf{3. PACF}\\ ordre $\hat p$};
    \node[box, below=6mm of pacf] (yw) {\textbf{4. Yule--Walker}\\ poids $\Phi_i$};
    \node[box, left=8mm of yw] (glr) {\textbf{5. Variance long terme}\\ Student corrigé};
    \node[box, left=8mm of glr] (res) {\textbf{6. Box--Pierce résidus}\\ tout absorbé ?};
    \draw[fl] (bp) -- (st);
    \draw[fl] (st) -- (pacf);
    \draw[fl] (pacf) -- (yw);
    \draw[fl] (yw) -- (glr);
    \draw[fl] (glr) -- (res);
  \end{tikzpicture}
  \end{center}
\end{frame}

% ======================================================================
%  15. BOX-PIERCE
% ======================================================================
\begin{frame}{1. Box--Pierce : les mois sont-ils indépendants ?}
  \[
    S_{\mathrm{BP}} = T\sum_{h=1}^{6}\hat\rho(h)^2
    \;\underset{H_0}{\sim}\; \chi^2(6)
  \]
  \vspace{3mm}
  \begin{itemize}\setlength{\itemsep}{4mm}
    \item sous $H_0$, chaque $\hat\rho(h)$ $=$ bruit d'ordre $1/\sqrt T$ :
      $\sqrt T\,\hat\rho(h)\to\mathcal N(0,1)$
    \item \textbf{6} décalages $=$ durée maximale des trades
      {\footnotesize (plafond vérifié en fin de chaîne)}
  \end{itemize}
  \vspace{5mm}
  \begin{block}{Chez nous}
    La moitié des séries autocorrélées $\to$ correction.\\
    Les autres (dont \texttt{oracle}, \texttt{rsi\_strict}) : verdict conservé
    tel quel.
  \end{block}
\end{frame}

% ======================================================================
%  16. DF + KPSS
% ======================================================================
\begin{frame}{2. Stationnarité : deux tests aux $H_0$ opposées}
  \begin{columns}[T]
    \begin{column}{0.52\textwidth}
      \small
      \textbf{Dickey--Fuller} \hfill {\footnotesize $H_0 : \phi=1$}
      \[
        x_t = \phi\,x_{t-1}+\eta_t
        \qquad
        t_{\mathrm{DF}} = \frac{\hat\phi - 1}{\widehat{\mathrm{se}}(\hat\phi)}
      \]
      \begin{itemize}\setlength{\itemsep}{1mm}
        \item $|\phi|<1$ : les chocs s'éteignent
        \item $\phi=1$ : ils s'accumulent (dérive)
        \item forme de Student, \alert{pas sa loi} :\\
          seuil $\approx-2{,}9$ (loi tabulée DF)
      \end{itemize}
    \end{column}
    \begin{column}{0.46\textwidth}
      \small
      \textbf{KPSS} \hfill {\footnotesize $H_0$ : stationnaire}
      \[
        S_t=\textstyle\sum_{s\le t}(x_s-\bar x)
      \]
      \begin{itemize}\setlength{\itemsep}{1mm}
        \item cumuls bornés $=$ stable
        \item cumuls qui s'emballent $=$ dérive
        \item couvre l'angle mort de DF
      \end{itemize}
    \end{column}
  \end{columns}
  \vspace{5mm}
  \begin{block}{Concordance exigée}
    DF \textbf{rejette} \emph{et} KPSS \textbf{ne rejette pas} $\to$ stationnaire.
    Sinon : « ? ».\\
    Chez nous : concordance partout.
  \end{block}
\end{frame}

% ======================================================================
%  17. AR(p) : PACF -> YULE-WALKER -> VARIANCE LT
% ======================================================================
\begin{frame}{3--5. Modéliser la mémoire, corriger la variance}
  \textbf{Ordre} --- PACF : impact \emph{propre} de chaque décalage
  \hfill {\footnotesize $\sqrt T\,\hat\tau(h)\sim\mathcal N(0,1)$, plafond 6}
  \vspace{3mm}

  \textbf{Poids} --- Yule--Walker : on multiplie le modèle AR par $x_{t-h}$,
  espérance :
  \[
    \rho(h) = \Phi_1\rho(h-1)+\cdots+\Phi_{\hat p}\,\rho(h-\hat p)
    \qquad
    \sigma^2_\eta = \gamma(0)\bigl(1-\textstyle\sum_i\Phi_i\rho(i)\bigr)
  \]
  \vspace{1mm}

  \textbf{Correction} --- toutes les autocovariances, pas seulement la
  diagonale :
  \[
    \gamma_{\mathrm{lr}} = \frac{\sigma^2_\eta}{\bigl(1-\sum_i \Phi_i\bigr)^2}
    \qquad
    \widehat{\mathrm{Var}}(\bar x)\approx\frac{\gamma_{\mathrm{lr}}}{T}
    \qquad
    \mathrm{DEFF}_t = \frac{\gamma_{\mathrm{lr}}}{\gamma(0)}
  \]
  \vspace{3mm}
  \begin{itemize}\setlength{\itemsep}{2mm}
    \item $\Phi=0$ : on retrouve $s^2/T$ --- l'analogue temporel du DEFF
    \item chez nous : $\hat p \le 3$, corrections jusqu'à $\times 6{,}4$
      $\to$ diapo suivante
  \end{itemize}
\end{frame}

% ======================================================================
%  18. RESULTATS SERIE MENSUELLE
% ======================================================================
\begin{frame}{6. Validation et résultats de l'acte III}
  \begin{center}
  \footnotesize
  \begin{tabular}{@{}lrrrrrc@{}}
    \toprule
    Stratégie & $T$ & $\hat\rho(1)$ & $\hat p$ & $\mathrm{DEFF}_t$ &
      $p$ corrigée & $p$ BP résidus \\
    \midrule
    \texttt{oracle}        & 193 & $0{,}18$ & 0 & $1{,}0$ & $9{,}4\times10^{-85}$ & $0{,}16$ \\
    \midrule
    \texttt{ma\_crossover} & 188 & $0{,}56$ & 3 & $\alert{6{,}4}$ & $0{,}94$ & $0{,}14$ \\
    \texttt{dt\_top}       & 197 & $0{,}38$ & 1 & $2{,}2$ & $0{,}97$ & $0{,}69$ \\
    \texttt{rsi\_classic}  & 195 & $0{,}23$ & 1 & $1{,}6$ & $0{,}08$ & $0{,}93$ \\
    \texttt{rsi\_strict}   & 138 & $0{,}03$ & 0 & $1{,}0$ & $0{,}26$ & $0{,}28$ \\
    \bottomrule
  \end{tabular}
  \end{center}
  \vspace{4mm}
  \begin{itemize}\setlength{\itemsep}{3mm}
    \item résidus : aucun rejet $\to$ correction légitime
    \item variance sous-estimée jusqu'à $\times\alert{6{,}4}$\ldots\
      \textbf{aucun verdict ne bascule}
  \end{itemize}
  \vspace{2mm}
  \begin{block}{}
    Fin de l'escalade : il ne reste que \structure{le témoin}.
  \end{block}
\end{frame}

% ======================================================================
%  19. EFFICIENCE
% ======================================================================
\begin{frame}{Aucune ne passe : est-ce normal ? Le marché répond}
  Rendements quotidiens du marché, 4\,132 jours :
  \vspace{2mm}
  \begin{center}
  \small
  \begin{tabular}{@{}lr@{}}
    \toprule
    Box--Pierce & $S_{\mathrm{BP}} = 93{,}7$, \quad $p \approx 0$ \\
    $\hat\rho(1)$ & $-0{,}08$ \\
    Part du lendemain expliquée & $0{,}6\,\%$ \\
    Gain de la stratégie « rebond » & $0{,}007\,\%$/jour, \alert{$<$ frais} \\
    \bottomrule
  \end{tabular}
  \end{center}
  \vspace{5mm}
  \begin{itemize}\setlength{\itemsep}{4mm}
    \item significatif, mais rien d'\textbf{exploitable} : efficience
      \emph{faible} (Fama)
    \item même leçon qu'à l'acte I : \structure{significativité $\neq$ taille
      d'effet}
  \end{itemize}
\end{frame}

% ======================================================================
%  20. PROTOCOLE RECAPITULE
% ======================================================================
\begin{frame}{Le livrable : un protocole en six étages}
  \begin{center}
  \small
  \begin{tabular}{@{}lll@{}}
    \toprule
    Étage & Menace traitée & Outil \\
    \midrule
    1. \edge{} apparié & dérive du marché & Monte Carlo apparié \\
    2. Student $+$ TCL & aléa d'échantillonnage & test de référence \\
    3. Bonferroni & multiplicité (11 tests) & risque global \\
    4. DEFF $+$ bootstrap & dépendance intra-titre & sondages, rééchant. \\
    5. Deux portes & pseudo-réplication & agrégation équipondérée \\
    6. Série AR($p$) & autocorrélation résiduelle & variance de long terme \\
    \bottomrule
  \end{tabular}
  \end{center}
  \vspace{5mm}
  \begin{itemize}\setlength{\itemsep}{3mm}
    \item réutilisable sur \textbf{tout signal futur} de l'agent
    \item condition invérifiable $\Rightarrow$ « ? » : s'abstenir et signaler
    \item converge avec les standards de l'économétrie financière
      {\footnotesize (Fama--MacBeth, Newey--West)}
  \end{itemize}
\end{frame}

% ======================================================================
%  21. LIMITES & PERSPECTIVES
% ======================================================================
\begin{frame}{Limites assumées, suites prévues}
  \textbf{Limites}
  \begin{itemize}\setlength{\itemsep}{2mm}
    \item résultats \emph{in-sample} $\to$ validation hors échantillon,
      période après période
    \item biais du survivant ; frais simplifiés
    \item corrections empilées : garantie exacte du seuil perdue si les trois
      canaux étaient forts \emph{en même temps} {\footnotesize (signalé ---
      absent ici)}
  \end{itemize}
  \vspace{5mm}
  \textbf{Perspectives}
  \begin{itemize}\setlength{\itemsep}{2mm}
    \item créer nos propres stratégies (apprentissage / test séparés)
    \item garde-fou d'effectif ; \edge{} modélisé selon le contexte
    \item signaux tournés vers l'avenir ; avance--retard (Granger)
  \end{itemize}
\end{frame}

% ======================================================================
%  22. CONCLUSION
% ======================================================================
\begin{frame}{Conclusion}
  \vspace{2mm}
  \begin{itemize}\setlength{\itemsep}{5mm}
    \item Le test de référence déclarait \textbf{deux} gagnantes ; la seule
      vérification de l'\textbf{indépendance} a renversé les deux verdicts
    \item Le \structure{témoin}, construit pour gagner, a franchi chaque étage
  \end{itemize}
  \vspace{5mm}
  \begin{block}{}
    Un outil de validation se juge à sa capacité d'attraper ses propres
    erreurs \emph{tout en reconnaissant le vrai} : la démonstration est faite.
  \end{block}
  \vspace{5mm}
  \begin{center}
    \structure{Merci de votre attention.}
  \end{center}
\end{frame}

% ======================================================================
%  ANNEXES — à dégainer pendant les questions
% ======================================================================
\appendix

\begin{frame}{Annexes}
  \small
  \begin{columns}[T]
    \begin{column}{0.5\textwidth}
      \begin{itemize}\setlength{\itemsep}{1mm}
        \item[A1.] Dérivation du DEFF
        \item[A2.] ICC : la formule de $n_0$
        \item[A3.] Shapiro--Wilk : la statistique $W$
        \item[A4.] Student : loi exacte, $n-1$ ddl
        \item[A5.] Bonferroni : exact vs Boole
        \item[A6.] DF : pourquoi pas la loi de Student ; AIC
      \end{itemize}
    \end{column}
    \begin{column}{0.5\textwidth}
      \begin{itemize}\setlength{\itemsep}{1mm}
        \item[A7.] $\hat\rho$ : bruit en $1/\sqrt T$
        \item[A8.] PACF : définition formelle
        \item[A9.] Yule--Walker : dérivation
        \item[A10.] Variance de long terme : preuve
        \item[A11.] Les deux oracles (v1 $\to$ v2)
        \item[A12.] Valeurs extrêmes : traitement
      \end{itemize}
    \end{column}
  \end{columns}
\end{frame}

% --- A1 ---------------------------------------------------------------
\begin{frame}{A1. Dérivation du DEFF}
  $n$ trades, $k$ titres de $m$ trades ($n = km$).
  Hypothèses : $\mathrm{Var}(x_i)=\sigma^2$ ;
  $\mathrm{Cov}(x_i,x_j)=\rho\sigma^2$ si même titre, $0$ sinon.
  \begin{align*}
    \mathrm{Var}\Bigl(\sum_i x_i\Bigr)
      &= \sum_i \mathrm{Var}(x_i) + \sum_i\sum_{j\neq i}\mathrm{Cov}(x_i,x_j)\\
      &= n\sigma^2 + n(m-1)\rho\sigma^2
      \qquad\text{\scriptsize (chaque trade : $m-1$ partenaires de grappe)}\\[1mm]
    \mathrm{Var}(\bar x)
      &= \frac{\sigma^2}{n}\,\bigl[\,1+(m-1)\rho\,\bigr]
      = \frac{\sigma^2}{n}\cdot\mathrm{DEFF}
  \end{align*}
  \begin{block}{Contrôles de cohérence}
    $\rho=0 \Rightarrow \mathrm{DEFF}=1$ (indépendance).\quad
    $\rho=1 \Rightarrow \mathrm{DEFF}=m \Rightarrow
      \mathrm{Var}(\bar x)=\sigma^2/k$ : il ne reste que $k$ informations
      (une par titre).\quad
    $m=1 \Rightarrow \mathrm{DEFF}=1$.\\
    Grappes inégales : $m \to \bar m$.
  \end{block}
\end{frame}

% --- A2 ---------------------------------------------------------------
\begin{frame}{A2. ICC : le terme $n_0$}
  \[
    \hat\rho = \frac{\mathrm{MSB}-\mathrm{MSW}}
                    {\mathrm{MSB}+(n_0-1)\mathrm{MSW}}
    \qquad
    n_0 = \frac{1}{k-1}\Bigl(N - \frac{\sum_i n_i^2}{N}\Bigr)
  \]
  \begin{itemize}\setlength{\itemsep}{2mm}
    \item $k$ titres, $N$ trades, $n_i$ trades du titre $i$
    \item $n_0$ = taille moyenne des grappes \textbf{ajustée} de leur
      inégalité : le terme $\sum n_i^2/N$ pénalise le déséquilibre
      (50 trades sur un titre, 3 sur un autre)
    \item grappes égales de taille $n$ : $n_0 = n$ ; un trade par titre :
      $n_0=1$, le terme intra disparaît
    \item rôle : plus les grappes sont grosses, plus il faut de contraste
      inter-titres pour un même $\hat\rho$ --- normalise la mesure entre 0 et 1
  \end{itemize}
\end{frame}

% --- A3 ---------------------------------------------------------------
\begin{frame}{A3. Shapiro--Wilk : la statistique $W$}
  \[
    W = \frac{b^2}{\mathrm{SCE}}, \qquad W \le 1
  \]
  \begin{itemize}\setlength{\itemsep}{2mm}
    \item $\mathrm{SCE} = \sum_i (x_i - \bar x)^2$ : dispersion classique
    \item $b^2$ : dispersion \emph{attendue si les valeurs triées suivaient
      une normale} (combinaison des statistiques d'ordre et de leurs
      espérances sous normalité)
    \item les deux ne coïncident que sous normalité $\to$ $W \simeq 1$ ;
      c'est la version chiffrée de la droite de Henry (QQ-plot)
    \item $H_0$ ($W\simeq1$) : normalité --- p-value faible $\Rightarrow$
      \textbf{rejet de la normalité} (sens inverse de Student !)
  \end{itemize}
  \begin{block}{Pourquoi il sur-rejette à grand $n$}
    Sous $H_0$, la distribution de $W$ se resserre vers 1 quand $n$ croît :
    le moindre écart devient significatif. Chez nous, $W$ de $0{,}50$ à
    $0{,}95$ --- \emph{tous} rejetés. D'où le recours au TCL.
  \end{block}
\end{frame}

% --- A4 ---------------------------------------------------------------
\begin{frame}{A4. Student : loi exacte et degrés de liberté}
  Si l'on connaissait $\sigma$ : $(\bar e - \mu)/(\sigma/\sqrt n)
  \sim \mathcal N(0,1)$ exactement (données normales).\\[2mm]
  $\sigma$ inconnu $\to$ remplacé par $s$, \emph{estimé sur le même
  échantillon} :
  \[
    t = \frac{\bar e-\mu}{s/\sqrt n}
    \;=\; \frac{\mathcal N(0,1)}{\sqrt{\chi^2_{n-1}/(n-1)}}
    \;\sim\; t_{n-1}
  \]
  \begin{itemize}\setlength{\itemsep}{2mm}
    \item rapport de deux variables aléatoires $\to$ queues plus épaisses
      que la normale : l'incertitude sur $s$ est absorbée par la loi
    \item $n-1$ : un degré consommé pour estimer $\bar e$ --- le même $n-1$
      que la correction de Bessel dans
      $s^2=\frac{1}{n-1}\sum(e_i-\bar e)^2$
    \item dès $n \gtrsim 100$ : Student $\approx$ normale (chez nous,
      $n \ge 978$)
  \end{itemize}
\end{frame}

% --- A5 ---------------------------------------------------------------
\begin{frame}{A5. Bonferroni : calcul exact vs borne de Boole}
  \textbf{Exact (tests indépendants), par le complémentaire :}
  \[
    \mathrm P(\text{au moins un FP})
    = 1 - \mathrm P(\text{aucun FP})
    = 1 - (1-\alpha)^{11}
    = 1 - 0{,}95^{11} \approx 0{,}43
  \]
  {\small On multiplie $0{,}95$ par lui-même (puissance) --- l'addition
  $11\times0{,}05=0{,}55$ ne serait qu'un majorant (et dépasserait 1 avec
  plus de tests).}
  \vspace{2mm}

  \textbf{Boole (sans hypothèse d'indépendance) :}
  \[
    \mathrm P\Bigl(\bigcup_i A_i\Bigr) \le \sum_i \mathrm P(A_i)
    \quad\Longrightarrow\quad
    \sum_i \alpha' = 0{,}05 \;\text{ garantit le risque global}
  \]
  \begin{itemize}
    \item $\alpha$ est conditionnel \emph{par définition} :
      $\alpha = \mathrm P(\text{rejeter } H_0 \mid H_0 \text{ vraie})$ ---
      le calcul se place dans le pire cas ($H_0$ partout)
    \item Vérification : $1-(1-0{,}0045)^{11}\approx0{,}049<0{,}05$
      (légèrement conservateur : possibles faux négatifs)
  \end{itemize}
\end{frame}

% --- A6 ---------------------------------------------------------------
\begin{frame}{A6. DF : pourquoi pas la loi de Student ; l'AIC}
  \textbf{La régression} : couples $(x_{t-1}, x_t)$, moindres carrés :
  \[
    \hat\phi = \frac{\sum x_{t-1}x_t}{\sum x_{t-1}^2}
    \qquad
    \widehat{\mathrm{se}}(\hat\phi)=\frac{\hat\sigma_\eta}{\sqrt{\sum x_{t-1}^2}}
  \]
  {\scriptsize En pratique on régresse la différence :
  $\Delta x_t = \delta x_{t-1}+\eta_t$, $\delta=\phi-1$, test $\delta=0$
  contre $\delta<0$ --- équivalent.}
  \vspace{1mm}

  \textbf{Pourquoi pas Student} : la loi du $t$ de régression suppose un
  régresseur \emph{stationnaire}. Sous $H_0$ ($\phi=1$), $x_{t-1}$ est une
  marche aléatoire, sa variance croît avec $t$ : les conditions s'effondrent.
  Loi tabulée par Dickey \& Fuller, décalée vers les négatifs :
  rejet à 5\,\% vers $-2{,}9$, contre $-1{,}65$ pour une normale.
  \vspace{1mm}

  \textbf{Version augmentée} : on ajoute $\Delta x_{t-1},\ldots,\Delta x_{t-k}$
  pour blanchir les résidus ; $k$ choisi par
  \[
    \mathrm{AIC}_k = 2\ln\hat\sigma_k + \frac{k}{T}
  \]
  {\small ajustement contre complexité (chez nous : $k$ de 0 à 9). Ces retards
  sont un \emph{correctif interne} au test --- rien à voir avec l'ordre
  $\hat p$ du modèle AR, qui vient de la PACF.}
\end{frame}

% --- A7 ---------------------------------------------------------------
\begin{frame}{A7. Pourquoi $\hat\rho(h)$ est un bruit en $1/\sqrt T$ sous $H_0$}
  \[
    \hat\rho(h) = \frac{\sum_{t=1}^{T-h}(x_t-\bar x)(x_{t+h}-\bar x)}
                       {\sum_{t=1}^{T}(x_t-\bar x)^2}
  \]
  \begin{itemize}\setlength{\itemsep}{2mm}
    \item sous $H_0$ (indépendance), chaque produit
      $(x_t-\bar x)(x_{t+h}-\bar x)$ a une \textbf{espérance nulle}
      (facteurs indépendants et centrés)
    \item mais sur $T$ fini, les produits positifs et négatifs ne se
      compensent qu'imparfaitement --- même phénomène qu'une moyenne
      d'échantillon jamais exactement égale à $\mu$
    \item TCL sur cette somme de termes non corrélés :
      $\mathrm{Var}(\hat\rho(h))\approx 1/T$, donc
      $\sqrt T\,\hat\rho(h)\to\mathcal N(0,1)$
    \item d'où $T\hat\rho(h)^2 \sim \chi^2(1)$, et par indépendance
      asymptotique entre décalages :
      $S_{\mathrm{BP}} = T\sum_{h=1}^{6}\hat\rho(h)^2 \sim \chi^2(6)$
  \end{itemize}
  \begin{block}{Ordre de grandeur}
    $T=195$ mois $\Rightarrow$ bruit typique $1/\sqrt{195}\approx0{,}07$ :
    des $\hat\rho$ de $\pm0{,}07$ sont attendus \emph{sans} dépendance.
  \end{block}
\end{frame}

% --- A8 ---------------------------------------------------------------
\begin{frame}{A8. PACF : définition formelle}
  \textbf{Le problème} : l'effet de chaîne. AR(1) de poids $0{,}6$ :
  $\rho(2)=0{,}36\neq0$ alors que $t-2$ n'a \emph{aucun} effet direct
  (le lien transite par $t-1$).
  \vspace{1mm}

  \textbf{Définition} : $\hat\tau(h)$ = corrélation entre $x_t$ et $x_{t-h}$
  \emph{une fois retirée l'influence linéaire des intermédiaires} :
  \begin{enumerate}
    \item régresser $x_t$ sur $(x_{t-1},\ldots,x_{t-h+1})$ $\to$ résidu $u_t$
    \item régresser $x_{t-h}$ sur les mêmes $\to$ résidu $v_t$
    \item $\hat\tau(h) = \mathrm{corr}(u_t, v_t)$
  \end{enumerate}
  {\small Équivalence : $\hat\tau(h)$ = dernier coefficient d'un AR($h$)
  ajusté sur la série.}
  \begin{block}{Pourquoi elle identifie l'ordre}
    Pour un AR($p$) : $\tau(h)=0$ pour tout $h>p$ (coupure nette), alors que
    $\rho(h)$ décroît sans s'annuler. Décision : sous $H_0$,
    $\sqrt T\hat\tau(h)\sim\mathcal N(0,1)$ $\to$ bande $\pm1{,}96/\sqrt T$ ;
    on garde le dernier décalage significatif (plafond 6, intermédiaires
    conservés --- surestimer la mémoire plutôt que la tronquer).
  \end{block}
\end{frame}

% --- A9 ---------------------------------------------------------------
\begin{frame}{A9. Yule--Walker : la dérivation}
  Série stationnaire \textbf{centrée} : $\gamma(h) = \mathrm{Cov}(x_t,x_{t-h})
  = \mathbb E[x_t x_{t-h}]$ --- l'espérance d'un produit de variables
  centrées \emph{est} leur covariance.\\[2mm]
  Modèle : $x_t = \Phi_1 x_{t-1}+\cdots+\Phi_p x_{t-p} + \eta_t$.
  On multiplie par $x_{t-h}$ ($h\ge1$), on prend l'espérance :
  \[
    \underbrace{\mathbb E[x_t x_{t-h}]}_{\gamma(h)}
    = \sum_i \Phi_i\,\gamma(h-i)
    + \underbrace{\mathbb E[\eta_t x_{t-h}]}_{=\,0\ (\eta_t \perp \text{passé})}
    \;\Longrightarrow\;
    \rho(h) = \sum_i \Phi_i\,\rho(h-i)
  \]
  {\small $p$ équations ($h=1,\ldots,p$), $p$ inconnues ; $\rho(0)=1$,
  $\rho(-k)=\rho(k)$. Cas $p=1$ : $\Phi_1=\rho(1)$.}
  \begin{block}{Décalage nul $\to$ le bruit}
    Avec $x_t$ lui-même, le bruit ne disparaît plus
    ($\mathbb E[\eta_t x_t]=\sigma^2_\eta$) :
    \[
      \gamma(0) = \sum_i\Phi_i\gamma(i) + \sigma^2_\eta
      \;\Longrightarrow\;
      \sigma^2_\eta = \gamma(0)\Bigl(1-\sum_i\Phi_i\rho(i)\Bigr)
    \]
    = variance totale moins la part expliquée par le passé.
  \end{block}
\end{frame}

% --- A10 --------------------------------------------------------------
\begin{frame}{A10. Variance de long terme : la preuve en trois pas}
  \textbf{1. Pourquoi $s^2/T$ est faux :}
  \[
    \mathrm{Var}(\bar x) = \frac{1}{T^2}\sum_t\sum_s \mathrm{Cov}(x_t,x_s)
    \;\approx\; \frac{1}{T}\sum_{h=-\infty}^{+\infty}\gamma(h)
    = \frac{\gamma_{\mathrm{lr}}}{T}
  \]
  {\small $s^2/T$ ne garde que la diagonale $\gamma(0)$ et jette tous les
  $\gamma(h)$ croisés.}
  \vspace{1mm}

  \textbf{2. Pour un AR($p$)} --- effet cumulé d'un choc : $+1$ aujourd'hui,
  $\sum\Phi$ demain, $(\sum\Phi)^2$ après\ldots\ son poids total :
  \[
    1+\Bigl(\sum\Phi\Bigr)+\Bigl(\sum\Phi\Bigr)^2+\cdots
    = \frac{1}{1-\sum\Phi}
    \;\Longrightarrow\;
    \gamma_{\mathrm{lr}} = \frac{\sigma^2_\eta}{(1-\sum_i\Phi_i)^2}
  \]
  \textbf{3. Vérification exacte sur l'AR(1)} :
  $\gamma(h)=\phi^{|h|}\gamma(0)$, $\gamma(0)=\sigma^2_\eta/(1-\phi^2)$ :
  \[
    \sum_h \gamma(h) = \gamma(0)\,\frac{1+\phi}{1-\phi}
    = \frac{\sigma^2_\eta}{(1-\phi)^2} \;\checkmark
  \]
  {\small Cohérence : $\Phi=0 \Rightarrow \gamma_{\mathrm{lr}}=\gamma(0)$,
  on retrouve $s^2/T$ ; $\sum\Phi \to 1$ : la correction explose.}
\end{frame}

% --- A11 --------------------------------------------------------------
\begin{frame}{A11. Les deux oracles (v1 $\to$ v2)}
  \textbf{Oracle v1} --- seuil $+25\,\%$ sur 30 jours :
  \begin{itemize}\setlength{\itemsep}{0.5mm}
    \item ne tradait que les rebonds extrêmes (2009, 2020) $\to$
      $t \approx 120$, série mensuelle « tassée », presque plate hors bouffées
    \item or DF a besoin de variabilité dans $x_{t-1}$ ($\hat\phi$ et se reposent sur
      $\sum x_{t-1}^2$) : erreur-type énorme, \textbf{aucune puissance} $\to$ DF ne rejette
      pas, KPSS dit stationnaire $\to$ divergence, verdict « ? »
  \end{itemize}
  \begin{block}{L'erreur de conception}
    Un témoin trop parfait produit une série dégénérée, \emph{hors du domaine
    de puissance} des outils qu'il doit calibrer. Un protocole à un seul test
    l'aurait écarté \textbf{en silence} ; la règle de concordance a rendu
    l'anomalie visible.
  \end{block}
  \textbf{Oracle v2} --- seuil $+5\,\%$, $\sim$1 an de repos entre entrées :
  \begin{itemize}\setlength{\itemsep}{0.5mm}
    \item entrées réparties sur tous les régimes, volume comparable aux
      vraies stratégies : témoin \textbf{réaliste}, avantage toujours net
    \item franchit les six étages --- et le témoin aura servi deux fois :
      contrôler le faux négatif \emph{et} réviser le protocole
  \end{itemize}
\end{frame}

% --- A12 --------------------------------------------------------------
\begin{frame}{A12. Valeurs extrêmes : gardées, et voici pourquoi}
  \textbf{Distinction} : valeur \emph{aberrante} = erreur de mesure (à
  corriger) $\neq$ valeur \emph{extrême} = vraie réalisation d'une queue
  lourde (de l'information).\\
  {\small Ici : pas d'aberrantes (cours réels, rendements bornés) mais des
  extrêmes massives --- visibles dans le $W$ de Shapiro ($0{,}50$ à $0{,}95$).}
  \vspace{1mm}
  \begin{block}{Ni supprimées, ni winsorisées --- quatre protections}
    \begin{enumerate}
      \item le test porte sur la \textbf{moyenne} : un extrême gonfle aussi
        $s$ au dénominateur (effet conservateur)
      \item le \textbf{bootstrap} les intègre sans hypothèse de forme
      \item l'\textbf{agrégation} leur retire leur poids : un méga-trade ne
        vote que dans sa grappe
      \item un avantage porté par quelques extrêmes est \emph{exactement} ce
        que les deux portes détectent (cas \texttt{rsi\_strict})
    \end{enumerate}
  \end{block}
  \textbf{Pourquoi pas la médiane ?} Le gain total $= n\times$ moyenne :
  une stratégie à médiane négative peut être rentable par quelques grands
  gains --- la médiane la déclarerait perdante à tort. La moyenne est le seul
  paramètre qui agrège les extrêmes \emph{avec} le reste.
\end{frame}

\end{document}
